% part: intuitionistic-logic
% chapter: semantics
% section: introduction

\documentclass[../../../include/open-logic-section]{subfiles}

\begin{document}

\olfileid[hi]{int}{sem}{int}

\olsection{परिचय}

अर्थविज्ञान के बिना किसी तर्कशास्त्र का संतोषजनक वर्णन नहीं होता, और
अंतःप्रज्ञावादी तर्कशास्त्र भी इसका अपवाद नहीं है। शास्त्रीय तर्कशास्त्र के
लिए !!{valuation}ों पर आधारित अर्थविज्ञान विहित है, जबकि अंतःप्रज्ञावादी
तर्कशास्त्र के लिए अनेक प्रतिस्पर्धी अर्थविज्ञान हैं। उनमें से कोई भी इस अर्थ
में पूर्णतः संतोषजनक नहीं है कि वह संयोजकों के अर्थों का अंतःप्रज्ञावादी रूप
से स्वीकार्य विवरण देता हो।

मोडल तर्कशास्त्रों के अर्थविज्ञान के समान, !!{relational model}ों पर
आधारित अर्थविज्ञान शायद सबसे लोकप्रिय है। इसमें !!{propositional variable} जगतों
को नियत किए जाते हैं और ये जगत किसी अभिगम्यता संबंध से जुड़े होते हैं। वह
संबंध सदैव आंशिक क्रम होता है; अर्थात वह स्वपरावर्ती, प्रतिसममित और
संक्रमणशील है।

सहज रूप से इन जगतों को ज्ञान-अवस्थाएँ या ``साक्ष्यगत स्थितियाँ'' समझा जा
सकता है। कोई अवस्था~$w'$,~$w$ से तभी अभिगम्य है जब हमारी समस्त जानकारी के
अनुसार $w'$ ज्ञान की एक संभव (भावी) अवस्था हो; अर्थात वह~$w$ पर ज्ञात बातों
के संगत हो। कोई प्रतिज्ञप्ति एक बार ज्ञात हो जाए तो फिर अज्ञात नहीं बन सकती;
अर्थात जब भी~$w$ पर $!A$ ज्ञात हो और~$Rww'$, तब~$w'$ पर भी $!A$ ज्ञात है।
इसलिए ``ज्ञान'' अभिगम्यता संबंध के सापेक्ष एकदिष्ट है।

यदि ज्ञानात्मक तर्कशास्त्र की तरह ``$!A$ ज्ञात है'' को ``सभी ज्ञानात्मक
विकल्पों में सत्य'' परिभाषित करें, तो~$w$ पर $!A \land !B$ ज्ञात है यदि सभी
ज्ञानात्मक विकल्पों में $!A$ और~$!B$ दोनों ज्ञात हों। किंतु ज्ञान एकदिष्ट
और $R$ स्वपरावर्ती है, इसलिए इसका अर्थ है कि~$w$ पर $!A \land !B$ तभी ज्ञात
है जब~$w$ पर $!A$ और $!B$ दोनों ज्ञात हों। इसी कारण~$w$ पर $!A \lor !B$
तभी ज्ञात है जब उनमें से कम-से-कम एक ज्ञात हो। अतः $\land$ और $\lor$ के
लिए संयोजकों की सत्यता-शर्तें शास्त्रीय तर्कशास्त्र की शर्तों से मिलती हैं।

परंतु सशर्त की सत्यता-शर्तें शास्त्रीय तर्कशास्त्र से भिन्न हैं।~$w$ पर
$!A \lif !B$ तभी ज्ञात है जब $Rww'$ वाला कोई ऐसा~$w'$ न हो जहाँ $!A$ तो
ज्ञात हो पर $!B$ ज्ञात न हो। यह इस शर्त के समान नहीं है कि~$w$ पर $!A$
अज्ञात हो या $!B$ ज्ञात हो। क्योंकि यदि~$w$ पर न $!A$ ज्ञात हो न $!B$, तो
कोई ऐसी भावी ज्ञानात्मक अवस्था~$w'$ हो सकती है जहाँ $Rww'$ और जहाँ $!A$
तो ज्ञात हो जाए, पर $!B$ ज्ञात न हो।

हम $\lnot !A$ तभी जानते हैं जब ऐसी कोई संभव भावी ज्ञानात्मक अवस्था न हो
जिसमें हम~$!A$ को जानते हों। यहाँ विचार यह है कि यदि $!A$ ज्ञेय होता, तो
किसी संभव भावी ज्ञानात्मक अवस्था में~$!A$ ज्ञात हो जाता। चूँकि हम $\lfalse$
को नहीं जान सकते, उस भावी ज्ञानात्मक अवस्था में हम $!A$ को जानते पर
~$\lfalse$ को नहीं जानते।

इस व्याख्या में बहिष्कृत मध्य का सिद्धांत विफल होता है। कुछ ऐसे $!A$ हैं
जिन्हें हम अभी नहीं जानते, पर भविष्य में जान सकते हैं। ऐसे !!{formula}~$!A$
के लिए $!A$ और~$\lnot !A$ दोनों अज्ञात हैं; अतः $!A \lor \lnot !A$ ज्ञात
नहीं है। किंतु हम, उदाहरणार्थ, $\lnot(!A \land \lnot !A)$ जानते हैं। क्योंकि
ऐसी कोई भावी अवस्था संभव नहीं है जिसमें हम $!A$ और $\lnot !A$ दोनों को
जानते हों; और यह हम इस बात से स्वतंत्र रूप से जानते हैं कि हमें~$!A$ या
$\lnot !A$ ज्ञात है या नहीं।

!!{relational model} अंतःप्रज्ञावादी तर्कशास्त्र का एकमात्र उपलब्ध
अर्थविज्ञान नहीं हैं। सांस्थितिक अर्थविज्ञान एक अन्य विकल्प है: इसमें
प्रतिज्ञप्तियों की व्याख्या किसी सांस्थितिक समष्टि के विवृत समुच्चयों के रूप
में और संयोजकों की व्याख्या इन समुच्चयों पर संक्रियाओं के रूप में की जाती
है (मसलन $\land$ सर्वनिष्ठ के अनुरूप है)।

\end{document}
